\documentclass[10pt,a4paper]{article}
\pdfinfoomitdate=1
\pdftrailerid{}
\pdfsuppressptexinfo=15

\usepackage[T1]{fontenc}
\usepackage[utf8]{inputenc}
\usepackage{lmodern}
\usepackage{microtype}
\usepackage[a4paper,margin=22mm]{geometry}
\usepackage{booktabs,tabularx,array}
\usepackage{enumitem}
\usepackage{xcolor}
\usepackage{hyperref}
\usepackage{xurl}
\usepackage{fancyhdr}

\definecolor{navy}{HTML}{17324D}
\definecolor{teal}{HTML}{1B728C}
\definecolor{soft}{HTML}{F3F6F8}
\definecolor{warn}{HTML}{8A4B08}
\definecolor{greenish}{HTML}{2E6B57}
\hypersetup{
  colorlinks=true,
  linkcolor=navy,
  urlcolor=teal,
  pdftitle={Siamese Embedding Compression Research Program v0.3.1 Status Addendum},
  pdfauthor={Guillaume Harbonnier},
  pdfsubject={Accepted v0.6 semantic hardening, N7/N8 transport hardening and Phase A construct-validity transition}
}
\setlength{\parindent}{0pt}
\setlength{\parskip}{5pt}
\setlist[itemize,enumerate]{nosep,leftmargin=*}
\pagestyle{fancy}
\fancyhf{}
\fancyhead[L]{Siamese Embedding Compression Research Program}
\fancyhead[R]{v0.3.1 status addendum}
\fancyfoot[C]{\thepage}
\renewcommand{\headrulewidth}{0.4pt}

\newcommand{\code}[1]{\texttt{#1}}

\title{\textbf{Siamese Embedding Compression Research Program}\\[4pt]
\large Status Addendum: semantic ACCEPT to construct-validity review}
\author{Guillaume Harbonnier\\
\small Independent exploratory research\\
\small \url{https://github.com/gharbonnier78/siamese-embedding-compression-lab}}
\date{Version 0.3.1 --- 16 September 2026}

\begin{document}
\maketitle
\thispagestyle{empty}

\begin{center}
\fcolorbox{warn}{soft}{\parbox{0.92\linewidth}{
\textbf{Scope.} This addendum updates the status of the v0.3 monograph without replacing it. It records the accepted
v0.6 pre-execution semantic review, the new N7/N8 review-transport findings, and the transition to a separate Phase A
construct-validity review. It opens no protected Study 1B SCREEN, qualification TEST, real-route performance or
representation geometry, launches no S4N3, and does not authorize canonical Phase A execution.
}}
\end{center}

\section{What changed since v0.3}

The v0.3 monograph described N4, N5, R1/N2 and N6 as prospectively hardened but still awaiting focused independent
review. That review is now complete on exact head
\code{c7d03cdad3bdf460bc6261508f2dc416504c7e4b}.

The independent verdict was:

\begin{quote}
\code{VERDICT: ACCEPT}\\
\code{N4\_CLOSED: yes}\\
\code{N5\_CLOSED: yes}\\
\code{R1\_N2\_PORTABLE\_CHOICE\_OVERLAP\_CLOSED: yes}\\
\code{N6\_PACKAGE\_RULE\_SATISFIED: yes}\\
\code{CANONICAL\_EXECUTION\_CURRENTLY\_ADMISSIBLE: no}\\
\code{CONSTRUCT\_VALIDITY\_REVIEWED: no}
\end{quote}

The reviewer independently reported 20/20 package SHA-256 and byte-length checks, 16/16 Git blob SHA-1 recomputations,
a PASS from the bundled verifier, 11/11 focused v0.6 contract tests, direct lock/addendum agreement re-derivation, and
no protected-scope leakage in the reviewed human-facing material.

The reviewed v0.6 lock and addendum are not rewritten merely to encode acceptance. The append-only review verdict and
Chronicle entry carry that decision while preserving the exact reviewed bytes.

\section{Meaning of the accepted semantic hardening}

\subsection*{N4: possession is not creation influence}

A genuinely pre-existing external source does not lose D2 eligibility merely because the attester later receives,
reviews or cites it. Production, commissioning, requesting, specifying, directing or materially framing evidence
requires D1. D1 remains the completeness backstop for every lock-defining choice and ambiguity fails closed.

\subsection*{N5: only temporal ordering is mechanically checkable}

The accepted wording no longer invites the entire D1/D2 mechanism to be called structurally objective. Only the temporal
ordering against the immutable repository-public bound is mechanically checkable. External/pre-existing classification,
absence of attester influence and completeness remain provenance/attestation claims.

\subsection*{R1/N2: informational overlap is choice-specific}

Prior evidence is classified against the lock-defining choice it may inform. Hardware mismatch alone cannot discharge
relevance for portable backend/BLAS, threading/affinity or measurement-method choices. Measurement-method evidence may
transfer even when hardware class, vector dimension or search family differ. If informative prior evidence has already
shaped a choice, the remedy is disclosure plus a new prospective design/execution identity and independent review, not
silent reinterpretation.

\subsection*{N6: exact bytes are reviewable}

The accepted package transported exact normative bytes and allowed the reviewer to recompute Git blob identity locally.
N6 is therefore satisfied for that reviewed package.

\section{New transport findings N7 and N8}

The same reviewer identified two new findings. They are transport-only and are not conditions of the v0.6 ACCEPT.

\begin{description}
\item[N7.] The package manifest exposed \code{git\_blob\_sha1\_expected} but did not explicitly declare that its source
was the repository object at the immutable reviewed head. Future packages must state that provenance explicitly.
The authored \code{git\_blob\_sha1\_recomputed} column should disappear so recomputation remains reviewer-side.

\item[N8.] The reviewer prompt required verification of the immutable PR binding and five exact-head bounded workflow
successes, but those objects were not transported. The next package must carry the binding snapshot and exact-head
workflow evidence, with direct GitHub/API locators retained for optional independent re-fetch.
\end{description}

The distinction is important: N7/N8 improve \emph{demonstrability and review autonomy}. They do not change the accepted
scientific or engineering semantics of N4/N5/R1-N2/N6.

\section{Why construct validity is now the next gate}

The governance question is no longer only ``is the benchmark protected against outcome-informed design?'' The next
question is:

\begin{quote}
\textbf{If Phase A were executed exactly as specified, would its measurements validly support the bounded engineering
claims the benchmark is designed to make?}
\end{quote}

This is a construct-validity question. It can and should be answered before target hardware or protected outcomes are
opened.

The dedicated review must challenge at least the following:

\begin{itemize}
\item whether 512D versus 128D is isolated as the intervention under otherwise paired execution conditions;
\item whether exact dense top-1 fixed-work search is a valid construct for the intended brute-force 1:N engineering question;
\item whether deterministic L2-normalized dense synthetic vectors are adequate for bounded memory/compute/latency/energy
characterization, and where numerical distribution could still matter;
\item whether G0 and G2 are valid sensitivity scenarios without being presented as event or deployment facts;
\item whether latency, throughput, RAM, queueing, external power and joules-per-identification endpoints match the claims;
\item whether external power measurement and off-device load generation provide the intended device-level energy construct;
\item whether counterbalancing, restarts, raw-sample retention and no-runtime-tuning rules control material confounding;
\item whether external-validity limits are explicit enough to prevent a bound platform/workload result from becoming a
universal edge-system claim;
\item whether characterization is kept separate from pass/fail qualification while operational thresholds remain null;
\item whether engineering evidence remains firewalled from biometric non-inferiority or superiority claims.
\end{itemize}

The formal request is stored in:\\
\path{protocol/reviews/STUDY1B_PHASE_A_CONSTRUCT_VALIDITY_REVIEW_REQUEST_V0_1_2026-09-16.md}.

\section{Current state}

\begin{tabularx}{\linewidth}{>{\bfseries}l X}
\toprule
Item & State \\
\midrule
Study 0 & corrected bounded negative closure \\
S4N1 / S4N2 & \code{CLOSED\_NEGATIVE} / \code{CLOSED\_NEGATIVE} \\
S4N3 & not launched \\
N4 / N5 / R1-N2 / N6 & accepted closed for reviewed v0.6 package \\
N7 / N8 & open review-transport hardening for next package \\
Construct validity & not reviewed \\
Target hardware / meter / load generator & not materialized \\
Execution-specific environment lock & not materialized \\
Canonical Phase A & not admissible; not executed \\
Protected SCREEN / TEST / performance / geometry & sealed \\
\bottomrule
\end{tabularx}

\section{Exact next admissible sequence}

\begin{enumerate}
\item implement N7/N8 in the construct-validity review-package transport;
\item synchronize human-readable documentation and freeze an exact review head;
\item require the five bounded technical-assurance workflows to succeed on that exact head;
\item create an immutable binding for that exact construct-validity review head;
\item build the N7/N8-hardened construct-validity package including binding and workflow evidence;
\item obtain independent construct-validity review;
\item only after construct-validity acceptance may the project materialize and freeze target hardware, external power
meter, off-device load generator, configuration-choice provenance, generator identity and an execution-specific lock;
\item rerun bounded assurance on the exact materialized execution head before any canonical Phase A execution becomes eligible.
\end{enumerate}

No step in this sequence authorizes opening protected Study 1B biometric outcomes.

\section{Authoritative artifacts added by this transition}

\begin{itemize}
\item \path{protocol/reviews/PR50_C7D03CDA_V06_N6_FOCUSED_REREVIEW_VERDICT_2026-09-16.md}
\item \path{protocol/chronicle/STUDY1B_PR50_C7D03CDA_V06_ACCEPT_2026-09-16.yaml}
\item \path{protocol/chronicle/STUDY1B_N7_N8_REVIEW_TRANSPORT_HARDENING_OPEN_2026-09-16.yaml}
\item \path{protocol/reviews/STUDY1B_PHASE_A_CONSTRUCT_VALIDITY_REVIEW_REQUEST_V0_1_2026-09-16.md}
\end{itemize}

Pinned reusable method authority:\\
\url{https://github.com/gharbonnier78/scientific-research-harness/blob/3b109adcdd9a8cba4df029d3803ee0e5cb5bdf98/HARNESS.md}

\vfill
\begin{center}
\small\textit{This addendum is explanatory and status-synchronizing. Frozen protocols, independent review verdicts and
append-only Chronicle entries remain authoritative.}
\end{center}

\end{document}
